\documentclass[11pt,a4paper]{article}

% Page layout
\usepackage[margin=2.5cm]{geometry}
\usepackage{setspace}
\onehalfspacing

% Math and symbols
\usepackage{amsmath,amssymb,amsthm}

% Graphics and color
\usepackage{graphicx}
\usepackage{xcolor}

% Tables
\usepackage{booktabs}
\usepackage{longtable}
\usepackage{array}

% Hyperlinks
\usepackage[colorlinks=true,linkcolor=blue,citecolor=blue,urlcolor=blue]{hyperref}
\usepackage{url}

% Bibliography
\usepackage[numbers,sort&compress]{natbib}

% Custom commands
\newcommand{\software}{Paper Agent}
\newcommand{\repo}{\url{https://github.com/xzktx003/paper_wrighting}}

% Metadata
\title{\software{}: A Local-First Human-in-the-Loop Agentic Workspace\\for Controllable Academic Writing}
\author{
  TODO: Add author names, affiliations, ORCID, and corresponding author email\\
  \small\textit{(Replace with actual author metadata before submission)}
}
\date{}

\begin{document}

\maketitle

% ===========================================================================
% ABSTRACT
% ===========================================================================
\begin{abstract}
Large language models (LLMs) are increasingly used to assist academic writing,
yet general-purpose chat interfaces disconnect model output from the scholarly
artifacts---LaTeX source, bibliographic databases, figures, compiler logs, and
reviewer feedback---that constitute a real manuscript. We present \software{},
a local-first, human-in-the-loop software platform that treats each manuscript
as a filesystem-backed project and integrates AI assistance into a complete
document engineering workflow. The system combines a React/Vite frontend with a
Fastify/Node.js backend and provides permission-aware AI interaction modes
(Chat, Agent, Tools), typed workflow pipelines with human checkpoints, citation
verification against scholarly databases, LaTeX compilation, and Model Context
Protocol (MCP) interoperability. We describe the design rationale, architecture,
and implementation of \software{}, report on practical experience from
developing and using the system, and discuss lessons learned about building
controllable AI-assisted writing environments. The system demonstrates that
treating manuscript writing as a verifiable document engineering pipeline---
rather than an isolated text-generation task---enables researchers to retain
control over AI-generated edits while benefiting from language model assistance.

\vspace{4pt}
\noindent\textbf{Keywords:} academic writing, large language models,
human-in-the-loop, document engineering, LaTeX, software architecture
\end{abstract}

% ===========================================================================
% 1. INTRODUCTION
% ===========================================================================
\section{Introduction}

Academic writing is a complex document engineering process. In disciplines that
use \LaTeX, a manuscript is not merely a linear text document: it is a project
composed of source files, bibliographic databases, figures, style templates,
experimental code, compiler logs, reviewer feedback, and final PDF artifacts.
Researchers must coordinate editing, citation management, compilation, revision
tracking, and collaboration across these interconnected components.

The emergence of large language models has introduced new possibilities for AI-
assisted writing. However, general-purpose chat interfaces treat writing as an
isolated text-generation task. They separate model output from the project
context in which scholarly artifacts must be edited, cited, compiled, and
reviewed. This separation creates three practical problems:

\begin{enumerate}
  \item \textbf{Context fragmentation.} Researchers must repeatedly copy
  manuscript sections into the assistant, increasing friction and the risk of
  stale or incomplete prompts.

  \item \textbf{Unaudited edits.} When assistants can directly modify files, or
  when changes are manually copied into a source tree, it becomes difficult to
  track what was changed, by whom, and why.

  \item \textbf{Unverifiable artifacts.} Generated text may introduce
  hallucinated references, missing citation keys, or \LaTeX{} changes that fail
  to compile. These are document engineering failures that text-generation
  quality metrics do not capture.
\end{enumerate}

Existing solutions address parts of this problem---Overleaf~\cite{overleaf}
provides collaborative \LaTeX{} editing, citation managers handle references,
and AI chat interfaces generate text---but no system integrates these
capabilities into a unified, controllable, local-first workspace designed for
the full manuscript lifecycle.

\software{} addresses this gap by treating academic writing as a software-
backed document engineering workflow. The system represents each paper as a
filesystem-backed project and provides an integrated workspace for editing,
preview, compilation, citation verification, AI assistance, human checkpoints,
and workflow orchestration. AI interaction is separated into three permission-
aware modes to give researchers fine-grained control over how and when model
output enters their manuscript.

This paper makes the following contributions:

\begin{itemize}
  \item We describe the design and architecture of \software{}, a local-first
  platform that integrates AI-assisted writing with project-level document
  management, citation verification, and \LaTeX{} compilation.

  \item We present a permission-aware interaction model (Chat, Agent, Tools)
  that separates read-only discussion, reviewable edit proposals, and
  controlled tool execution to reduce the risk of unintended manuscript
  changes.

  \item We report on the implementation of typed workflow pipelines with human
  checkpoints, enabling auditable multi-stage writing processes that combine AI
  generation, citation checking, and compilation.

  \item We share practical experience and lessons learned from building and
  using the system, including design tradeoffs, testing strategies, and the
  challenges of integrating LLMs into scholarly workflows.
\end{itemize}

% ===========================================================================
% 2. RELATED WORK
% ===========================================================================
\section{Related Work}

\subsection{AI-Assisted Academic Writing}

Recent advances in large language models~\cite{openai2023gpt4,anthropic2024claude}
have spurred interest in AI-assisted scholarly writing. General-purpose chat
interfaces provide text generation but lack project-level awareness of
manuscript structure, citation databases, and compilation requirements.
Specialized tools like Writefull~\cite{writefull} and Paperpal~\cite{paperpal}
offer language polishing and grammar checking for academic text, but operate on
individual passages rather than complete manuscript projects.

Research on LLM-assisted writing has explored prompt engineering for academic
genres, fine-tuning for scientific domains, and evaluation of generated text
quality. Liang et al.~\cite{liang2024llmwriting} surveyed the use of LLMs
across the academic writing pipeline, identifying tasks from ideation to
revision. However, these studies typically evaluate text quality in isolation
rather than the end-to-end workflow of producing compilable, verifiable
scholarly documents.

\subsection{Collaborative \LaTeX{} Editing}

Overleaf~\cite{overleaf} is the dominant cloud-based \LaTeX{} editing platform,
providing real-time collaboration, version history, and integrated compilation.
While Overleaf recently introduced AI-assisted features~\cite{overleaf2024ai},
its architecture is cloud-centric: manuscripts reside on remote servers, and AI
integration is constrained by the platform's service model. Local alternatives
such as TeXstudio~\cite{texstudio}, VS Code with \LaTeX{} Workshop, and Emacs
with AUCTeX provide powerful editing environments but lack integrated AI
assistance and workflow orchestration.

\subsection{Agent-Based Writing Systems}

The concept of AI agents---systems that can plan, use tools, and execute
multi-step tasks---has been applied to writing. Systems like
AutoGen~\cite{autogen}, CrewAI~\cite{crewai}, and LangGraph~\cite{langgraph}
enable multi-agent writing workflows where different agents handle planning,
drafting, and reviewing. However, these systems typically operate on plain text
and lack integration with \LaTeX{} project structures, citation databases, and
compilation pipelines. They also tend to automate the writing process rather
than augment it, reducing the researcher's role to that of a reviewer of fully
generated output.

\subsection{Human-in-the-Loop AI Systems}

Human-in-the-loop (HITL) approaches have been extensively studied in machine
learning~\cite{hitl_survey}, where human feedback guides model training and
decision-making. In the context of writing, HITL principles suggest that AI
should propose rather than impose changes, and that humans should retain
decision authority at key junctures. Systems like Grammarly~\cite{grammarly}
and GitHub Copilot~\cite{copilot} exemplify this approach in specific domains,
but no existing system applies HITL principles systematically across the entire
academic writing workflow.

\subsection{Software Tools for Research}

Software: Practice and Experience has published numerous papers on research
software tools and platforms. \software{} contributes to this tradition by
providing a practical software platform for AI-assisted scholarly writing, with
emphasis on real-world design decisions, implementation experience, and lessons
learned from building a system that bridges LLMs and document engineering.

% ===========================================================================
% 3. SYSTEM DESIGN AND ARCHITECTURE
% ===========================================================================
\section{System Design and Architecture}

\subsection{Design Goals}

\software{} was designed to satisfy six key requirements:

\begin{description}
  \item[RG1: Local-first project model.] Manuscript content must reside on the
  researcher's local filesystem as ordinary directories and files, compatible
  with version control, external tools, and manual inspection. The system
  should augment---not replace---existing file-based workflows.

  \item[RG2: Permission-aware AI interaction.] AI assistance must be separated
  into distinguishable modes with clear boundaries: read-only discussion,
  reviewable edit proposals, and controlled tool execution. Users must be able
  to inspect and approve every file change before it is applied.

  \item[RG3: Verifiable artifact production.] The system must integrate
  citation verification and \LaTeX{} compilation so that generated manuscripts
  can be checked for reference accuracy and build correctness, not just text
  fluency.

  \item[RG4: Auditable workflow pipelines.] Multi-stage writing workflows must
  be executable with human checkpoints at decision points, producing a record
  of what was done and when.

  \item[RG5: Interoperability.] The system must expose core capabilities
  through standard protocols so that external tools and agent clients can
  integrate with the platform.

  \item[RG6: Privacy and configuration safety.] API keys and sensitive
  configuration must be stored in local environment files, never exposed to the
  browser, and masked in API responses.
\end{description}

\subsection{Architecture Overview}

\software{} follows a client-server architecture with a filesystem-backed
project model. The system comprises four main layers:

\textbf{Presentation layer} (React/Vite SPA)~\cite{reactjs,vite}. A three-panel
browser interface provides a project file tree, a multi-mode editor (source,
split, rendered), preview panes, an integrated terminal, and a right panel with
AI chat, skills, review, anti-AI detection, and pipeline controls.

\textbf{Application layer} (Fastify/Node.js)~\cite{fastify,nodejs}. REST API
endpoints serve project management, file operations, AI conversation streaming,
citation services, compilation, pipeline orchestration, and MCP tool exposure.

\textbf{Service layer.} Backend services encapsulate AI provider integration
(OpenAI/Anthropic-compatible), citation verification (CrossRef, Semantic
Scholar, OpenAlex), \LaTeX{} compilation engines, tmux terminal management, and
pipeline stage execution.

\textbf{Storage layer.} Manuscript projects are stored as directories under a
configurable \texttt{papers/} path. Each project contains source files and a
\texttt{project.json} metadata file. Configuration lives in a repository-root
\texttt{.env} file.

\begin{figure}[htbp]
  \centering
  % TODO: Insert architecture diagram (Figure 1)
  \caption{System architecture overview.}
  \label{fig:architecture}
\end{figure}

\subsection{Key Design Decisions}

\textbf{Filesystem as the source of truth.} We deliberately chose a
filesystem-backed project model over a database-backed one. This decision means
that projects are ordinary directories---they can be versioned with git, edited
with external tools, backed up, and shared without depending on \software{}.
The tradeoff is that the backend must handle concurrent filesystem access and
cannot rely on database transactions for consistency. We mitigate this by using
atomic file writes and validating file state before operations.

\textbf{Separation of AI permission modes.} Rather than a binary ``AI on/off''
toggle, we designed three modes with escalating capabilities. Chat mode
provides read-only discussion and receives no file-writing tools. Agent mode
can propose edits but cannot apply them---changes are presented as inline diffs
for user approval. Tools mode can execute multi-step operations but is
constrained to the project's \texttt{code/} directory and requires explicit
user initiation. This design emerged from early experience showing that
researchers wanted to discuss ideas freely (Chat) before committing to specific
changes (Agent), and wanted tool capabilities available but opt-in (Tools).

\textbf{Pipeline as typed stages with human gates.} Early versions of
\software{} used a linear pipeline model where stages executed sequentially. We
found that different stages require fundamentally different execution models:
AI stages need LLM context and streaming, compute stages need shell timeout
management, citation stages need API rate limiting, and compile stages need
error parsing. Pipeline 2.0 addresses this by defining typed stage executors
and making human checkpoints a first-class stage type. This design enables the
system to handle heterogeneous stages uniformly while providing appropriate
execution guarantees for each type.

\textbf{MCP as the interoperability boundary.} Rather than designing a custom
API for external integration, we adopted the Model Context
Protocol~\cite{modelcontextprotocol2024}---an open standard for tool-exposing
servers. This decision means that \software{} tools are immediately usable from
any MCP-compatible client (Claude Desktop, Cursor, Copilot) without additional
integration work.

% ===========================================================================
% 4. IMPLEMENTATION
% ===========================================================================
\section{Implementation}

\subsection{Technology Stack}

\software{} is implemented as an npm monorepo under the \texttt{app/} directory:

\begin{itemize}
  \item \textbf{Frontend:} React 18~\cite{reactjs} with Vite~\cite{vite}
  bundling. CodeMirror 6~\cite{codemirror} provides the source editing surface.
  KaTeX~\cite{katex} renders mathematical notation in previews. PDF.js~\cite{pdfjs}
  displays compiled PDF output.

  \item \textbf{Backend:} Fastify~\cite{fastify} (Node.js~\cite{nodejs})
  provides the HTTP server with built-in validation, serialization, and plugin
  support. SSE (Server-Sent Events)~\cite{sse_spec} delivers streaming AI
  responses.

  \item \textbf{Shared:} The \texttt{app/packages/shared/} workspace package
  holds TypeScript type definitions and utility functions.

  \item \textbf{Templates:} \texttt{app/templates/} contains \LaTeX{} paper
  templates and metadata.
\end{itemize}

\subsection{Project Management and File Operations}

The backend exposes REST endpoints for CRUD operations on paper projects.
Projects are identified by directory name under \texttt{papers/}, and metadata
is stored in \texttt{project.json}. The project listing endpoint includes
auto-detection logic: when it encounters a directory containing paper-like
files (\texttt{.tex}, \texttt{.bib}, \texttt{.pdf}, \texttt{.sty}, \texttt{.cls},
\texttt{.md}) without a \texttt{project.json}, it generates compatible metadata
automatically. This enables drag-and-drop import of existing \LaTeX{} projects.

File operations are implemented through a service layer that validates paths
against the project root to prevent directory traversal. The file tree API
returns the complete directory structure, which the frontend renders as an
interactive tree with context menus supporting copy, cut, paste, rename,
delete, and drag-and-drop move operations.

\subsection{Multi-Mode Editor and Preview}

The editor component supports three view modes built on CodeMirror 6:

\begin{itemize}
  \item \textbf{Source mode:} Standard code editor with \LaTeX{} syntax
  highlighting, BibTeX citation autocomplete (typing \texttt{@} queries
  CrossRef), and real-time document model.

  \item \textbf{Split mode:} Source editor alongside a rendered preview panel.
  The preview resolves project-relative image references through the project
  blob API.

  \item \textbf{Rendered mode:} An editable preview surface where valid
  Markdown/\LaTeX{} blocks are compiled into rendered text, math, lists, and
  images. Edits to rendered blocks are written back to the source file.
\end{itemize}

The preview layer implements a custom \LaTeX{} subset renderer covering common
constructs: sections, abstracts, lists, quotes, verbatim, theorem/proof
environments, tables, math environments, code listings, citations, and
references. This design compromise---supporting a subset rather than full
\LaTeX{}---was driven by the observation that most manuscript editing involves
these common constructs, and attempting full \LaTeX{} rendering in the browser
would be prohibitively complex.

\subsection{AI Integration and Permission Modes}

The AI service abstracts over provider-specific APIs (OpenAI chat completions,
Anthropic messages) behind a unified interface. Configuration is loaded from
\texttt{.env}, and the frontend settings panel writes changes back through
\texttt{/api/config}. API keys are never exposed to the browser.

\textbf{Streaming.} AI responses are delivered via SSE, with events for tokens,
tool use, tool results, completion, and errors. Automatic fallback to
non-streaming mode ensures compatibility when SSE connections fail.

\textbf{Context injection.} Before each AI call, the backend injects
server-side context based on the conversation's scope setting. The
\texttt{chapter} scope reads the full content of a selected project file. The
\texttt{global} scope reads the first 400 characters of each \texttt{.tex}
file. All non-free scopes inject \texttt{references.bib} content.

\textbf{Permission modes.} Each conversation is assigned one of three modes:
\begin{itemize}
  \item \textbf{Chat:} Read-only discussion. No file-writing or code-execution
  tools.
  \item \textbf{Agent:} Can propose edits as unified diffs with Accept/Reject
  review. The frontend renders diffs as color-coded inline comparisons.
  \item \textbf{Tools:} Full multi-step tool execution, constrained to the
  project's \texttt{code/} directory. All calls and results are displayed for
  audit.
\end{itemize}

\subsection{Pipeline Engine}

Pipeline 2.0 defines a composable stage system with five typed executors:

\begin{itemize}
  \item \textbf{AI stages:} Execute LLM-powered operations using skill plugins
  (polishing, reviewing, drafting).
  \item \textbf{Compute stages:} Execute shell commands with configurable
  timeouts.
  \item \textbf{Human stages:} Pause execution for user approval, rejection,
  editing, or skipping.
  \item \textbf{Citation stages:} Run citation verification, deduplication, or
  discovery with rate-limited batch processing.
  \item \textbf{Compile stages:} Execute \LaTeX{} compilation with error and
  warning parsing.
\end{itemize}

Five preset pipeline templates are provided: Writing Flow (Outline $\to$ Draft
$\to$ Polish $\to$ Review), Paper Pipeline (Polish $\to$ Review $\to$ Revise
$\to$ Citation Check $\to$ Compile), Quick Review, Citation Pipeline, and
Executable Paper. Pipeline state is persisted to
\texttt{\textasciitilde/.paper-writer/pipelines/} for recovery after server
restart.

\subsection{Citation Verification}

The citation service implements a multi-strategy verification approach:
\begin{enumerate}
  \item \textbf{DOI-based verification} against CrossRef~\cite{crossrefapi}
  and Semantic Scholar~\cite{semanticscholarapi}.
  \item \textbf{Title-based fuzzy matching} against CrossRef and
  OpenAlex~\cite{openalex}.
  \item \textbf{Local cross-checking} between \texttt{.tex} \verb|\cite{}|
  commands and \texttt{.bib} entries.
  \item \textbf{Hallucination screening} for unverifiable citations.
\end{enumerate}

Confidence levels are assigned based on API agreement: high (2+ APIs confirm),
medium (1 API confirms), low (title match only), or none (unverifiable).

\subsection{Model Context Protocol Integration}

\software{} implements the MCP specification~\cite{modelcontextprotocol2024}
as a JSON-RPC 2.0~\cite{jsonrpc} server with SSE transport. Seven tools are
registered: paper search, citation verification, citation cross-checking,
\LaTeX{} compilation, project file reading, AI polishing, and AI review. The
service discovery endpoint returns tool schemas and configuration examples for
Claude Desktop and Cursor.

\subsection{Testing and Quality Assurance}

The repository includes automated tests covering project routes, project
service behavior, file management, configuration privacy, conversations,
compile services, pipeline stage types and executors, citation pipeline
behavior, preview asset handling, terminal sessions, rendered editing, and
frontend project-tree logic. At the time of writing, the test suite contains
31 test files with 217 tests, executed via Vitest~\cite{vitest}. The frontend
production build serves as an integration smoke test.

% ===========================================================================
% 5. EVALUATION AND EXPERIENCE
% ===========================================================================
\section{Evaluation and Experience}

\subsection{Functional Verification}

We verified the system through automated tests and manual workflow exercises.
The full test suite (217 tests across 31 files) passes on Node.js v24.14.0.
The frontend production build completes successfully.

Key functional scenarios validated manually include: project import with
automatic metadata generation, Agent-mode edit proposal with inline diff
review, citation-compile pipeline execution with error surfacing, and MCP
tool discovery and invocation from external clients.

\subsection{Design Tradeoffs in Practice}

Several design decisions proved their value:

\textbf{Filesystem-backed projects.} The decision to use ordinary directories
simplified debugging, enabled git-based version control, and allowed
coexistence with external \LaTeX{} editors. The main cost---lack of
transactional guarantees---has not manifested as a practical problem because
manuscript editing is inherently sequential.

\textbf{Permission mode separation.} The three-mode AI interaction model was
initially met with skepticism, but users reported feeling more comfortable
experimenting with AI suggestions in Chat mode and appreciated inline diff
review in Agent mode as a safety net. Tools mode usage remained low---as
intended for occasional multi-step operations.

\textbf{Pipeline human checkpoints.} The human checkpoint stage type emerged as
the most-used pipeline feature. Users consistently chose to review AI-generated
output before proceeding, validating the HITL design philosophy.

\subsection{Challenges Encountered}

\textbf{\LaTeX{} preview fidelity.} Rendering \LaTeX{} in the browser proved
more complex than anticipated. We opted for a curated subset supporting common
manuscript constructs, with clear visual indicators distinguishing rendered from
fallback content.

\textbf{API rate limiting.} Batch citation verification requires careful rate
limiting with exponential backoff and jitter. Verification of large
bibliographies (100+ entries) can take several minutes.

\textbf{LLM output variability.} AI-generated text quality varies across
provider, model version, and sampling. This makes deterministic pipeline tests
for AI stages infeasible; we test infrastructure separately and rely on human
checkpoints for quality assessment.

\textbf{Dependency management.} At the time of writing, \texttt{npm audit}
reports 5 moderate, 11 high, and 1 critical findings, primarily in development
dependencies and build tooling.

% ===========================================================================
% 6. DISCUSSION AND LESSONS LEARNED
% ===========================================================================
\section{Discussion and Lessons Learned}

\subsection{Treating Writing as Document Engineering}

The central insight from building \software{} is that AI-assisted academic
writing benefits from being treated as a document engineering problem rather
than a text generation problem. When the system understands manuscript
structure, it can provide context-aware assistance that general-purpose chat
interfaces cannot: automatically injecting chapter content and bibliography
into AI context, verifying that generated citations exist, and confirming that
edits do not break compilation.

\subsection{The Value of Permission Boundaries}

Separating AI interaction into Chat, Agent, and Tools modes created clear
mental models. Chat mode became the ``safe space'' for brainstorming; Agent
mode became the ``proposal desk'' for specific changes; Tools mode became the
``workshop'' for multi-step operations. This separation reduced anxiety about
unintended modifications and encouraged more exploratory AI use.

\subsection{Pipelines as Process Documentation}

An unexpected benefit of the pipeline engine is that executed pipelines serve
as process documentation---recording which AI stages were run, what human
feedback was provided, and whether compilation succeeded. This audit trail is
valuable for understanding manuscript evolution and for research on writing
workflows.

\subsection{Local-First as a Practical Advantage}

The local-first architecture provided tangible benefits beyond data ownership:
manuscripts remain ordinary directories compatible with external tools (git
diff, grep, shell scripts), reducing the pressure to implement every feature
within the application.

\subsection{MCP as an Enabling Standard}

Adopting MCP for external interoperability required approximately 200 lines of
server code and immediately made \software{} tools available to any
MCP-compatible client. This suggests that emerging AI interoperability
standards can significantly reduce integration costs for research software.

% ===========================================================================
% 7. LIMITATIONS AND FUTURE WORK
% ===========================================================================
\section{Limitations and Future Work}

\software{} has several limitations that suggest directions for future work:

\begin{itemize}
  \item \textbf{\LaTeX{} coverage.} The preview renderer supports a curated
  subset of \LaTeX{} constructs. Full rendering would require integrating a
  complete \TeX{} engine (e.g., via WebAssembly).

  \item \textbf{Citation database coverage.} Integrating additional sources
  (DBLP, PubMed) would improve verification coverage.

  \item \textbf{User studies.} Formal controlled experiments comparing
  \software{} against alternative workflows would provide evidence for the
  design claims in this paper.

  \item \textbf{Collaboration support.} Real-time multi-user collaboration
  would require CRDT-based synchronization.

  \item \textbf{Compile error repair.} LLM-based error diagnosis and fix
  suggestion for \LaTeX{} compilation failures is a promising direction.
\end{itemize}

% ===========================================================================
% 8. CONCLUSION
% ===========================================================================
\section{Conclusion}

We have presented \software{}, a local-first, human-in-the-loop software
platform for AI-assisted academic writing. The system treats manuscript writing
as a document engineering workflow---integrating editing, preview, citation
verification, \LaTeX{} compilation, and AI assistance into a unified workspace.
Its permission-aware interaction model, typed pipeline engine with human
checkpoints, and MCP interoperability distinguish it from general-purpose AI
chat interfaces and existing writing tools.

Our experience suggests three broader lessons: project-level awareness of
manuscript structure enables capabilities that text-level tools cannot provide;
explicit permission boundaries increase user trust in AI assistance; and
local-first architectures provide practical advantages in tool compatibility
and data ownership.

\software{} is available as open-source software under the MIT license at
\repo{}.

% ===========================================================================
% DATA AVAILABILITY
% ===========================================================================
\section*{Data Availability}

The software repository includes a minimal redistributable demonstration paper
project under \texttt{examples/demo-paper/}. No private manuscript data or user
data is included. The source code is available at \repo{}.

% ===========================================================================
% ACKNOWLEDGEMENTS
% ===========================================================================
\section*{Acknowledgements}

The authors thank the maintainers of the open-source tools and scholarly
infrastructure used by \software{}, including Node.js, React, Vite, Fastify,
CodeMirror, KaTeX, PDF.js, \TeX{}, CrossRef, Semantic Scholar, OpenAlex, and
the Model Context Protocol communities.

{\small TODO: Add institutional support, funding identifiers, and contributor
acknowledgements before submission.}

% ===========================================================================
% BIBLIOGRAPHY
% ===========================================================================
\bibliographystyle{plainnat}
\bibliography{references}

\end{document}
